\documentclass[a4paper,11pt]{article}
\usepackage{fontspec}
\usepackage{polyglossia}
\setmainlanguage{vietnamese}
\setmainfont{DejaVu Serif}
\setsansfont{DejaVu Sans}
\setmonofont{DejaVu Sans Mono}
\usepackage[margin=2.2cm]{geometry}
\usepackage{titlesec}
\usepackage{enumitem}
\usepackage{longtable}
\usepackage{booktabs}
\usepackage{array}
\usepackage{xcolor}
\usepackage{hyperref}
\usepackage{parskip}
\usepackage{seqsplit}

\hypersetup{
  colorlinks=true,
  linkcolor=blue,
  urlcolor=blue
}

\titleformat{\section}{\Large\bfseries}{\thesection.}{0.6em}{}
\titleformat{\subsection}{\large\bfseries}{\thesubsection}{0.6em}{}

\newcolumntype{L}[1]{>{\raggedright\arraybackslash}p{#1}}

\title{\textbf{Project Scope \& Limitations}\\[4pt]
\large Dự án: Hệ thống Quản lý Du thuyền (Cruise Management System)\\
\normalsize Jira: SCRUM-9 -- Project Scope \& Limitations}
\author{}
\date{}

\begin{document}
\maketitle
\vspace{-1.5cm}
\tableofcontents
\vspace{0.5cm}

\section{Mục tiêu tài liệu}

Tài liệu này tổng hợp \textbf{Major Features} của toàn bộ 5 module trong hệ thống (Cruise/Itinerary, Activity/Excursion, Passenger/Booking, POS/Payment, Dashboard/Admin), đồng thời nêu rõ \textbf{Limitations \& Exclusions} -- những giới hạn và phạm vi mà hệ thống \emph{không} xử lý -- nhằm thống nhất phạm vi (scope) giữa các thành viên trong nhóm trước khi triển khai.

\section{Phạm vi dự án -- Major Features theo module}

\subsection{Module 1: Cruise / Itinerary (Quản lý chuyến đi \& lịch trình)}

\begin{itemize}[leftmargin=1.2cm]
  \item Quản lý thông tin tàu (Cruise/Ship): tên tàu, sức chứa, số tầng, số cabin.
  \item Quản lý chuyến đi (Voyage/Cruise Trip): ngày khởi hành, ngày kết thúc, điểm đi/đến.
  \item Quản lý lịch trình (Itinerary): danh sách điểm dừng (port), thời gian cập bến/rời bến từng điểm.
  \item Quản lý trạng thái chuyến đi: sắp diễn ra, đang diễn ra, đã hoàn thành, hủy.
  \item Xem lịch trình chi tiết theo từng ngày trong hành trình.
\end{itemize}

\subsection{Module 2: Activity / Excursion (Hoạt động \& Tour tham quan)}

\begin{itemize}[leftmargin=1.2cm]
  \item Quản lý danh mục hoạt động trên tàu (Activity): tên, mô tả, thời gian, địa điểm, số chỗ.
  \item Quản lý danh mục tour tham quan (Excursion): tên tour, điểm đến, giá, thời lượng, số chỗ.
  \item Quản lý lịch hoạt động/tour theo từng chuyến đi.
  \item Theo dõi số lượng đăng ký / số chỗ còn lại theo thời gian thực.
\end{itemize}

\subsection{Module 3: Passenger / Booking (Hành khách \& Đăng ký)}

\begin{itemize}[leftmargin=1.2cm]
  \item Quản lý thông tin hành khách (Passenger).
  \item Quản lý Booking (đặt chỗ chuyến đi).
  \item Quản lý Cabin (phòng ở trên tàu).
  \item Đăng ký hoạt động trước và trong chuyến đi (ActivityRegistration).
  \item Đăng ký tour tham quan (ExcursionRegistration).
  \item Check-in bằng QR Code, Cruise Card hoặc RFID/NFC (CheckIn).
  \item Xem lịch trình chuyến đi, dịch vụ miễn phí/có phí, chi tiêu phát sinh và hóa đơn tạm tính.
  \item Gửi phản hồi (feedback) sau hoạt động hoặc sau chuyến đi.
\end{itemize}

\subsection{Module 4: POS / Payment (Bán hàng \& Thanh toán)}

\begin{itemize}[leftmargin=1.2cm]
  \item Quản lý điểm bán hàng trên tàu (POS): nhà hàng, quầy bar, cửa hàng lưu niệm, spa...
  \item Ghi nhận giao dịch mua hàng/dịch vụ của hành khách trong chuyến đi.
  \item Quản lý phương thức thanh toán: Cruise Card, thẻ ngân hàng, ví điện tử.
  \item Tính hóa đơn tổng hợp (folio) theo từng hành khách/cabin, xuất hóa đơn cuối chuyến.
  \item Áp dụng khuyến mãi/giảm giá (nếu có) cho dịch vụ trên tàu.
\end{itemize}

\subsection{Module 5: Dashboard / Admin (Quản trị \& Báo cáo)}

\begin{itemize}[leftmargin=1.2cm]
  \item Quản lý tài khoản người dùng \& phân quyền (Admin, Nhân viên, Hành khách).
  \item Dashboard tổng quan: số chuyến đi, tỷ lệ lấp đầy cabin, doanh thu, số lượng đăng ký hoạt động/tour.
  \item Báo cáo thống kê theo chuyến đi, theo module, theo thời gian.
  \item Cấu hình danh mục dùng chung (loại cabin, loại hoạt động, điểm đến...).
  \item Nhật ký hoạt động hệ thống (audit log) ở mức cơ bản.
\end{itemize}

\section{Limitations \& Exclusions}

Các nội dung sau \textbf{nằm ngoài phạm vi} của hệ thống trong giai đoạn triển khai hiện tại:

\begin{itemize}[leftmargin=1.2cm]
  \item Không tích hợp trực tiếp với hệ thống đặt vé máy bay hoặc phương tiện di chuyển đến cảng khởi hành.
  \item Không xử lý nghiệp vụ hải quan, xuất nhập cảnh hoặc visa cho hành khách quốc tế.
  \item Không bao gồm hệ thống định vị/điều hướng tàu (navigation system) hay các hệ thống kỹ thuật vận hành tàu (engine, an toàn hàng hải...).
  \item Không tích hợp cổng thanh toán ngân hàng thật (chỉ mô phỏng/giả lập giao dịch thanh toán trong phạm vi đồ án).
  \item Không hỗ trợ đa ngôn ngữ (multi-language) ở phiên bản đầu tiên -- mặc định tiếng Việt.
  \item Không bao gồm ứng dụng di động (mobile app) riêng cho hành khách -- chỉ triển khai trên nền web.
  \item Không xử lý quản lý nhân sự thủy thủ đoàn (crew management) ngoài phạm vi tài khoản nhân viên vận hành hệ thống.
  \item Không đảm bảo hoạt động offline khi tàu mất kết nối Internet giữa hành trình (giả định luôn có kết nối).
\end{itemize}

\end{document}
